% TOP_PROBLEM: {"id": 10000036, "title": "Invariant finite-energy percolation with internal threshold one", "category_id": 19, "category": {"id": 19, "name": "probability", "display_name": "Probability", "description": "Problems involving probability theory, stochastic processes, and random structures.", "slug": "probability", "order_index": 19, "created_at": "2026-07-31T15:26:25.671Z"}, "status": "open", "problem_number": "AMR-099-0036", "set_id": 13}
% FIRST_POSTED: 2026-10-01
% ENTRY_KIND: statement_only
% UnsolvedMath ID 10000036; source catalogue code AMR-099-0036
% !TeX program = lualatex
% Definition and sources only; this file is not an LLM solution attempt.
\documentclass[11pt,a4paper]{article}
\usepackage[margin=25mm]{geometry}
\usepackage{fontspec}
\setmainfont{lmroman10-regular.otf}[BoldFont=lmroman10-bold.otf,
 ItalicFont=lmroman10-italic.otf,BoldItalicFont=lmroman10-bolditalic.otf]
\usepackage{amsmath,amssymb,mathtools}
\usepackage{unicode-math}
\setmathfont{latinmodern-math.otf}
\AtBeginDocument{\let\setminus\smallsetminus}
\usepackage{xurl,hyperref}
\hypersetup{colorlinks=true,urlcolor=blue}
\setcounter{secnumdepth}{0}
\setlength{\parindent}{0pt}
\setlength{\parskip}{5pt}
\setlength{\emergencystretch}{3em}
\begin{document}
\section*{Invariant finite-energy percolation with internal threshold one}\label{10000036}
{\small UnsolvedMath ID 10000036 · AMR-099-0036 · Probability · AMR Open Problem Lists}
\subsection{Definitions and mathematical statement}
For some integer $d\ge2$, let $E_d$ be the nearest-neighbor edge set of $\mathbb Z^d$. Seek a random spanning subgraph $X=(\mathbb Z^d,E_X)$ whose law is invariant under every automorphism of the lattice and satisfies the finite-energy condition
\[
0<\mathbb P(e\in E_X\mid E_X\setminus\{e\})<1
\quad\text{almost surely for every }e\in E_d.
\]
Here the conditioning records the states of all edges other than $e$. Require that $X$ have an infinite connected component almost surely.

Given one realization $X$, retain each of its edges independently with probability $q$; write $\mathbb P_q^X$ for this additional, conditional product law. Define its internal percolation threshold by
\[
p_c(X)=\inf\{q\in[0,1]:\mathbb P_q^X(\text{an infinite component exists})>0\}.
\]
Does there exist such an invariant finite-energy law with $p_c(X)=1$ almost surely?

Thus $X$ must percolate, but every independent thinning by a fixed positive deletion probability must destroy all infinite components. The second-stage Bernoulli randomness is independent of the randomness used to construct $X$. Finite energy refers to the original edge law; it is not a uniform lower bound on conditional opening probabilities, which would impose a stronger requirement. The question asks for existence in a lattice dimension and naturally also invites determining the dimensions in which it can occur. The threshold is defined on the entire possibly disconnected subgraph, so success in any one infinite component counts as percolation.
\subsection{Short English statement}
Can an invariant finite-energy lattice percolation percolate while every independent positive thinning destroys percolation?
\subsection{Sources}
\begin{itemize}
\item[S1] ULAM.AI and the UnsolvedMath Contributors, supplied problems.json, record 10000036 (AMR-099-0036), AMR Open Problem Lists. Catalogue identity and source transcription; CC BY 4.0.
\url{https://huggingface.co/datasets/ulamai/UnsolvedMath}.
\item[S2] Benjamini - Coarse Geometry and Randomness (Saint-Flour notes), Open problem 8.15, PDF page 61. Original source indexed by AMR-099-0036.
\url{https://www.wisdom.weizmann.ac.il/~itai/stflouraug24.pdf#page=61}.
\item[S3] I. Benjamini, Coarse Geometry and Randomness, primary Saint-Flour notes; the numbered question and its surrounding definitions.
\url{https://arquivo.pt/noFrame/replay/20201231041548id_/http://www.wisdom.weizmann.ac.il/~itai/stflouraug24.pdf}.
\item[S4] I. Benjamini and V. Tassion, Homogenization via sprinkling, Question 1, repeating the precise invariant finite-energy question.
\url{https://projecteuclid.org/journals/annales-de-linstitut-henri-poincare-probabilites-et-statistiques/volume-53/issue-2/Homogenization-via-sprinkling/10.1214/16-AIHP746.pdf}.
\end{itemize}
\end{document}
